\section{Methods}

\subsection{Dataset and Cohort}
We use the UK Biobank (UKB)~\cite{ukbiobank}, a prospective cohort of
500{,}000 UK adults aged 40--69 at recruitment.
We restrict to participants aged $\geq$\,60 at the end of follow-up or
death, yielding $n\,=\,124{,}158$ participants.
The binary outcome is all-cause mortality after age 60.
We apply a stratified 70/15/15 split: train\,86{,}911,
val\,18{,}624, test\,18{,}623 (event rate 26.2\% in all splits).
Features extracted prior to age 60 include: age at recruitment, sex, BMI,
smoking status, alcohol status, living alone, 21 blood biomarkers
(HbA1c, total cholesterol, HDL/LDL, CRP, cystatin~C, albumin, etc.),
11 binary comorbidity flags (T2 diabetes, CKD, depression, stroke, etc.),
and the full temporal history of ICD-10 diagnoses.

\subsection{Patient Data Representations}
\label{sec:representations}
The five methods differ in which modalities they use and how they
represent them (Fig.~\ref{fig:methods}):

\paragraph{Structured features (39-dim).}
Demographics, biomarkers, and comorbidity flags are concatenated into a
39-dimensional numeric vector used by the Baseline and as the raw-feature
supplement for embedding methods.

\paragraph{Disease-history prose.}
ICD events before age 60 are serialised chronologically:
\textit{``At age 44.0, patient was diagnosed with J45 vasomotor rhinitis.
At age 52.3, patient was diagnosed with F32 depressive episode.''}
Only disease events are included; no demographics or lab values.

\paragraph{Temporal ICD token stream.}
The same ICD events are formatted as a structured token sequence:
\textit{``44yr:J45 $\rightarrow$ 52yr:F32 $\rightarrow$ \ldots''}, preserving
the temporal ordering in a compact, model-native format.

\paragraph{Full clinical summary prose.}
All four modalities are combined into a single paragraph:
demographics, lab values with units, comorbidity flags, and disease
history. This is the only representation that includes numeric biomarker
values as text.

\subsection{Methods Compared}

\paragraph{Delphi (zero-shot).}
Delphi~\cite{delphi2024} is a 7B-parameter autoregressive LLM pre-trained
on tokenised EHR sequences.
We feed the temporal ICD token stream and extract the per-position
death-token logit.
A sigmoid hazard chain converts per-step logits to a survival curve:
$h_t = \sigma(\ell_t)$, $S(t) = \prod_{i \le t}(1 - h_i)$.
Risk at each horizon is $1 - S(t_{\text{horizon}})$.
No task-specific fine-tuning is applied.

\paragraph{Baseline CoxPH.}
The 39-dim structured feature vector is fitted directly with
$\ell_2$-penalised Cox proportional hazards (penaliser\,=\,0.1,
implemented in \texttt{lifelines}~\cite{lifelines}).

\paragraph{Disease Text Emb + CoxPH.}
The disease-history prose is encoded by frozen Qwen3-8B~\cite{qwen3}
in encoder mode (mean pooling of final hidden states) to a 4096-dim
vector, reduced to 128 dims via PCA fitted on the training set, then
concatenated with the 39 raw features (167-dim total) and fitted with CoxPH.

\paragraph{Trajectory Emb + CoxPH.}
Identical to Disease Text Emb + CoxPH, but the input is the temporal
ICD token stream instead of prose.
This isolates the effect of temporal ordering vs.\ prose formatting of
the same ICD events.

\paragraph{Clinical Summary Emb (no raw features).}
The full clinical summary prose is encoded by Qwen3-8B $\rightarrow$ PCA(128).
Crucially, \emph{no raw numeric features are concatenated}; the 128-dim
embedding is the sole input to CoxPH.
This tests whether the LLM embedding alone captures biomarker-level
discriminative structure.

\subsection{Evaluation Protocol}
We report three metrics on the held-out test set:
(i) Harrell C-index (concordance, higher is better)~\cite{harrellcindex};
(ii) mean time-dependent AUC (TD-AUC) averaged over horizons
    9.5, 15.4, and 19.7 years (higher is better);
(iii) Integrated Brier Score (IBS, lower is better), measuring calibration.
All metrics are computed with \texttt{lifelines} and \texttt{scikit-survival}.
No information leakage: PCA and CoxPH are fitted on training splits only;
val and test transforms use the training-fit PCA.

\begin{figure}[t]
  \centering
  \includegraphics[width=\textwidth]{figures/figure_methods}
  \caption{Methods overview. Patient data is separated into four modalities
           (demographics, biomarkers, comorbidity flags, temporal ICD history)
           and routed into five processing pipelines with distinct encoder
           architectures: autoregressive LLM (Delphi), no encoder (Baseline),
           and bidirectional Qwen3-8B encoder (embedding methods).}
  \label{fig:methods}
\end{figure}
